\section{INTRODUCTION}

Language-model assistance in engineering must connect an instruction to a change in geometry, and then to a manufacturing process whose outcome can be checked. A geometrically valid solid does not establish that the instruction was understood. A syntactically valid numerical control program does not establish that the intended material was removed. In machining, excess removal is irreversible, whereas residual stock can sometimes be removed by a subsequent operation. These distinctions require separate measures of semantic fidelity, execution success and geometric acceptance.

CAD datasets such as ABC \cite{Koch2019ABC} have enabled learning on engineering geometry. DeepCAD represents a solid as a sequence of modelling commands \cite{Wu2021DeepCAD}; Text2CAD pairs natural-language descriptions with command sequences \cite{Khan2024Text2CAD}. CAD-Coder trains a model to emit CadQuery scripts using supervised fine-tuning followed by reinforcement learning with a geometric reward \cite{Guan2025CADCoder}. Graph-CAD introduces an explicit assembly hierarchy and a constraint-decomposition graph as an intermediate representation \cite{GraphCAD2026}. Plan then Act splits generation between a planner that emits high-level intent and an actioner that produces low-level commands \cite{Guo2026PTA}. These studies ask how a model can be made to generate geometry correctly, and they measure reconstruction error, validity and geometric fidelity. Their progress is real: Plan then Act reports that direct low-level command generation by a general-purpose model fails 70.35 percent of the time, reduced to 0.85 percent by their two-stage design. A further line of work puts an agent directly in front of a computer-aided design kernel. SolidCraft \cite{SolidCraft2026} exposes a headless command engine over the Model Context Protocol, so a model can sketch, constrain, build and measure parametric solids. Its checks are geometric: volume, area, centre of mass and bounding box against analytic or FreeCAD-computed references. It emits no toolpath, no machine instructions and no removal volume, and it has no test for material that should have survived a cut.
We ask a complementary question. When the output actually removes material, what should the model be permitted to do, and how should the system decide whether the outcome is acceptable? The distinction matters because a valid command sequence is not necessarily a safe one. Work on language models and numerical control has so far evaluated or assisted the model's own output: comparing generated three-axis programs \cite{Sket2025}, checking tool and parameter advice against a digital twin \cite{Enis2026}, or combining program text with machine-interface screenshots for fault detection \cite{Pour2025}. Verification in these settings rests on human inspection, commercial computer-aided manufacturing software, or visual comparison. Planning and acting in physical environments has likewise been studied through reflection and retry loops \cite{Yao2023ReAct,Shinn2023Reflexion} and through standardised tool exposure \cite{Anthropic2024MCP}, where the model typically judges its own progress.

We study tool-mediated grounding in non-text engineering states. The \emph{spatial modality} consists of boundary-representation geometry and sampled height fields; the \emph{manufacturing modality} consists of operation graphs, toolpaths and simulated removal. The language models access these states through textual digests and typed calls, not through a newly trained multimodal encoder. The work concerns interaction with geometry rather than improved perceptual fusion or general multimodal intelligence. Deterministic evaluation makes a criterion reproducible, but does not make that criterion complete: a generator and its evaluator can agree while both miss the same error.

The architecture we built connects fifteen language models to a computer-aided design and manufacturing pipeline through typed tool contracts and an independently checkable verifier. Face-level machining features are recognised on boundary-representation adjacency graphs \cite{Colligan2022}, then passed through a gateway that exposes a closed set of operations. Region-masked toolpaths confine tool centres to positions that cannot reach material above the current cutting layer. A grid-based removal simulator turns the resulting program into gouging and residual volumes with attribution back to individual operations. A state machine, a blackboard, six separated roles, bounded repair and episodic memory close the loop. FreeCAD 1.1.4 serves as the reference geometry kernel \cite{FreeCAD2025}; the gateway backend is replaceable, but we report no cross-kernel claim.

The experiments ask three questions. First, how do contract calls and free-form scripts differ on the same Chinese work orders, and does the difference survive a reference construction independent of the contract implementation? Second, what do deterministic repair and memory add on ten anchor designs with two stock allowances? Third, where does agreement between a discrete generator and its verifier fail under an external geometric check? Fifteen-model contract evaluations probe remaining language-understanding errors, but are not a factorial experiment crossing every model with every action space. Direct program generation on ten anchor parts is a separate diagnostic, not the control arm of the twenty-work-order comparison.

The contribution is an evaluation procedure that follows errors across these interfaces. It separates export success from removed-material fidelity, audits reference geometry as well as model output, and tests the removal verifier using CAMotics and OpenCAMLib \cite{CAMotics2020,OpenCAMLib2023}. The external audit finds 316,398 flagged positions among 9,660,095 sampled tool centres on 130 parts, despite all 130 passing the internal criterion. Some violations exceed the 0.5 mm safety threshold. This disagreement limits the deployment claim, but makes the testbed useful for studying how apparently reliable interfaces can conceal shared errors.

Two negative results further constrain interpretation. A conventional facing operator does not implement the layer-aware machining required by these parts, and raising only the endpoints of model-written paths does not supply missing horizontal coverage. Neither observation establishes that a full manufacturing package or a model allowed to replan would fail. Our conclusions concern the evaluated interfaces and budgets; they do not isolate model scale as a cause or establish a universal advantage for closed action spaces.

The system generates its numerical control programs automatically; restricting the model's low-level output does not introduce a human programming step. We call this a relocation of autonomy: models propose semantic actions while deterministic components execute and check them. The restrictions provide an auditable division of responsibility, not a guarantee of physical safety. The preceding structured-generation literature already motivates intermediate representations; the additional contribution here is testing the reference and verification boundaries of a material-removal pipeline.

The remainder of the paper is organised as the journal requires. Materials and Methods defines the task, the acceptance criterion, the pipeline, the verification chains and the statistical treatment. Results and Discussion reports the recognition, toolpath, criterion, orchestration and model experiments, followed by the third-party audit of the verifier. Conclusion states what transfers, what does not, and what would have to change before any of this could be used on a shop floor.
